\chapter*{Conclusion générale}
\addcontentsline{toc}{chapter}{Conclusion générale}

Le projet Tastify a permis de concevoir et de réaliser une application web de gestion de restaurant couvrant plusieurs besoins opérationnels : gestion du menu, plan de salle, commandes, cuisine, stocks, ressources humaines, réservations, paiements, fidélité, avis clients et tableaux de bord.

Sur le plan technique, le projet met en œuvre une architecture moderne fondée sur Django REST Framework, React, MySQL, Redis, Celery et Docker Compose. L'intégration de WebSocket améliore la coordination entre la salle et la cuisine, tandis que les tâches asynchrones permettent de traiter certains calculs ou analyses sans bloquer l'expérience utilisateur.

Le projet intègre également une première couche d'intelligence applicative : analyse de sentiment des avis clients via Hugging Face, recommandations simples de plats et estimation des besoins à partir des ventes récentes. Ces fonctionnalités constituent une base solide qui pourra être enrichie dans de futures versions.

Les perspectives d'amélioration concernent principalement la personnalisation des recommandations, l'ajout d'un vrai modèle prédictif pour les stocks, la consolidation des tests de charge, l'amélioration de l'expérience mobile et l'éventuel déploiement sur une infrastructure de production.

En conclusion, Tastify représente une solution cohérente et évolutive pour la gestion numérique d'un restaurant. Le projet nous a permis de mobiliser des compétences en conception, développement backend, développement frontend, bases de données, architecture logicielle, sécurité et validation.
